\section{Conclusion}\label{sec:conclusion}

I have described an extension of \galpy\ that makes it run, compile, and
differentiate under \jax\ and \torch, on CPUs and GPUs, while leaving
its existing \numpy\ and \cc\ implementation byte-for-byte unchanged
and its interface as it was. Potentials, orbit integration, the
action--angle transformations and their inverses, and the distribution
functions, including those of tidal streams, accept the arrays of either
framework and return derivatives with respect to their inputs and
parameters. The existing test suite of \galpy, now 4,275 tests, passes
under \numpy, under both frameworks in eager mode, and under both with
compilation, with six tests remaining as known exceptions, and the
dedicated tests of the backend code add 8,585 more. Derivatives were
tested against finite differences whose convergence is demonstrated,
against the \cc\ variational equations, and against each other across
the two frameworks; where \galpy\ implements a numerical method of its
own, it was tested against \scipy\ and against arbitrary-precision
references. The work also found and fixed about fifty defects in the
\numpy\ and \cc\ code that had been there before.

What this makes possible is the gradient-based fitting of dynamical
models that \galpy\ already contains, with a likelihood whose derivative
with respect to every parameter---of the potential, of the DF, of the
orbits of the objects being modeled---costs a small multiple of the
likelihood itself, and the use of these models as components of larger
differentiable programs, such as neural networks that represent a
potential or a DF, or a simulation whose output is compared with data.
The speed-ups that a GPU provides, of one to three orders of magnitude
for large batches, are a second benefit, and one that depends on
keeping the computation on the device by compiling it as a whole.
Fitting a stream model to data with a Hamiltonian Monte Carlo sampler,
which the derivatives of \sectionname~\ref{sec:streams} allow, is the
obvious next application and will be presented separately.

Some things are not done. The torus-mapping code is not covered and will
not be. \astropy\ units cannot be used inside compiled programs. The
construction of the quasi-isothermal disk DF and of the stream-gap model
are \numpy\ computations, as are the variational equations and Lyapunov
exponents, which are features of the \cc\ integrators. Derivatives
through the \cc\ integrators and the \cc\ action--angle methods are with
respect to the initial conditions only and of first order. And a
backend is a moving target: the special-function fallbacks exist
because the frameworks' own versions are incomplete today, and the test
that checks this will say when they are no longer needed.

The code described here is part of \galpy's main branch from version
\flag{v1.13} on %TODO(JB): version
and is installed with \texttt{pip install galpy[jax]} or
\texttt{galpy[torch]}; the documentation describes the backends at
\\ \centerline{\url{https://docs.galpy.org}~.} All code used to produce
the figures in this paper, together with the working notes and logs of
the development, is available at
\\ \centerline{\url{https://github.com/jobovy/galpy-jax}~.}
